\chapter{Specification}\label{ch:b2:specification}

A regression answers the question its specification asks. Which variables enter, in what form, measured how, and how
many specifications were tried before one was reported: each choice changes the answer, sometimes predictably. This
chapter goes through the common choices. The chart shows the most predictable of the effects: a regressor measured with
noise gives a slope that shrinks toward zero by a factor the noise determines.

\begin{openingchart}{Estimated slope against the variance of the measurement error in the regressor, one simulated sample
  of 1,000 per point; true slope 2, variance of the true regressor 1. Curve: the limit $2/(1 + \sigma_e^2)$.}{Simulated
  (seeded); our computation}
\begin{tikzpicture}
\begin{axis}[omrecord, width=0.62\linewidth, height=0.4\linewidth, xmin=0, xmax=3, ymin=0, ymax=2.2,
    xlabel={Variance of the measurement error, $\sigma_e^2$}, ylabel={Estimated slope},
    legend style={at={(1.03,0.5)}, anchor=west, font=\scriptsize}]
  \addplot[omDef, only marks, mark=*, mark size=1.3pt] table[x=s2, y=slope]
    {figdata/bachelor-2/32-specification/out/attenuation.csv};
  \addlegendentry{estimate}
  \addplot[omThm, thick, domain=0:3, samples=100] {2/(1 + x)};
  \addlegendentry{$\lambda\beta$}
  \addplot[black!50, dashed, domain=0:3, samples=2] {2};
  \addlegendentry{true slope}
\end{axis}
\end{tikzpicture}
\end{openingchart}

% figdata: attenuation (s2 = 0: 1.994; 1: 0.990; 3: 0.507)
With no measurement error the estimate is 1.99; with an error as variable as the regressor itself it is 0.99, half the
true slope; with three times as much, 0.51. More data would not help: the points sit on the curve, not around the true
value.

\begin{recall}
From the Year 1 volume: logarithms and percentage changes, elasticities, $R^2$, tests and $p$-values. From
\cref{ch:b2:multiple-regression}: the \omterm{def:b2:multiple-regression:ovb}{omitted-variable bias} formula and \omterm{def:b2:multiple-regression:collinearity}{perfect collinearity}. From
\cref{ch:b2:regression-inference}: robust standard errors.
\end{recall}

\section{Logarithms and elasticities}\label{sec:b2:specification:logs}

\begin{definition}[Constant-elasticity regression]\label{def:b2:specification:log-log}
A \emph{constant-elasticity regression}\index{constant-elasticity regression} (log-log model) regresses $\ln y$ on $\ln
x$; its slope is the elasticity of $y$ with respect to $x$, the same at every value of $x$.
\end{definition}

\begin{definition}[Semi-elasticity]\label{def:b2:specification:semi-elasticity}
In a regression of $\ln y$ on $x$, the slope $\beta$ is a \emph{semi-elasticity}\index{semi-elasticity}: a unit change in
$x$ changes $y$ by approximately $100\beta$ per cent, exactly by $100(e^\beta - 1)$ per cent.
\end{definition}

The listing fits an Engel curve, food spending against income, to 235 household budgets in three ways.

\omcode{code/bachelor-2/32-specification/listings/engel.py}{1}{24}{An Engel curve for food in levels, in logs and as a
budget share.}\label{lst:b2:specification:engel}

\omoutput{code/bachelor-2/32-specification/outputs/engel.txt}{The log-log elasticity, the elasticity implied by the
linear fit at the means, and the slope of the budget share.}\label{out:b2:specification:engel}

% code: listings/engel.py; ledger: b2:engel-source
In logs, a 1 per cent rise in income goes with 0.856 per cent more food spending: food is a necessity, with an income
elasticity below one. The linear fit gives 0.485 francs of food per extra franc, an elasticity of 0.764 at the means: the
two forms answer the same question differently because the linear one forces the elasticity to change along the curve.
The budget share, 0.656 on average, falls by 0.086 for each unit of log income. Which form is right is an empirical
question (Exercise 7 lets the elasticity vary).

\section{Dummy variables and interactions}\label{sec:b2:specification:dummies}

\begin{definition}[Dummy variable]\label{def:b2:specification:dummy}
A \emph{dummy variable}\index{dummy variable} takes the value 1 if an observation belongs to a category and 0 otherwise;
its coefficient is the difference in the mean outcome between the category and the omitted one, other regressors fixed.
\end{definition}

\begin{definition}[Dummy-variable trap]\label{def:b2:specification:dummy-trap}
The \emph{dummy-variable trap}\index{dummy-variable trap} is the \omterm{def:b2:multiple-regression:collinearity}{perfect collinearity} created by including a constant and
a dummy for every category of a variable; one category must be left out (or the constant dropped).
\end{definition}

\begin{definition}[Interaction term]\label{def:b2:specification:interaction}
An \emph{interaction term}\index{interaction term} is the product of two regressors; with it the effect of one regressor
depends on the value of the other.
\end{definition}

With $\ln\mathit{food} = \beta_0 + \beta_1\ln\mathit{inc} + \beta_2\mathit{owner} + \beta_3\,\mathit{owner}\times\ln\mathit{inc} + u$,
the income elasticity is $\beta_1$ for tenants and $\beta_1 + \beta_3$ for owners. The listing applies it to Brenland's
household survey, with a dummy for each household size but the first.

\omcode{code/bachelor-2/32-specification/listings/dummies.py}{1}{23}{Size dummies, an owner dummy and an interaction in
Brenland's food regression.}\label{lst:b2:specification:dummies}

\omoutput{code/bachelor-2/32-specification/outputs/dummies.txt}{Household size absorbs a third of the raw elasticity;
owners' elasticity is slightly higher.}\label{out:b2:specification:dummies}

% code: listings/dummies.py
Without controls the elasticity is 0.288; with size dummies it is 0.196, because larger households have both more income
and more mouths to feed. A household of five spends 0.738 log points, $e^{0.738} - 1 = 109$ per cent, more on food than a
single person with the same income. Being an owner shifts nothing on its own, but owners' elasticity is 0.024 higher than
tenants'.

\section{Measurement error}\label{sec:b2:specification:me}

\begin{definition}[Classical measurement error]\label{def:b2:specification:classical-me}
A regressor has \emph{classical measurement error}\index{classical measurement error} if it is observed as $x = x^* + e$,
where the error $e$ has mean zero and is independent of the true value $x^*$ and of the regression error.
\end{definition}

\begin{definition}[Reliability ratio]\label{def:b2:specification:reliability}
The \emph{reliability ratio}\index{reliability ratio} of a measured variable is the share of its variance due to the true
value: $\lambda = \operatorname{Var}(x^*)/(\operatorname{Var}(x^*) + \operatorname{Var}(e))$.
\end{definition}

\begin{proposition}[Attenuation]\label{prop:b2:specification:attenuation}
If $y = \beta_0 + \beta x^* + u$ and $x = x^* + e$ with classical error, the least-squares slope of $y$ on $x$ converges in
probability to $\lambda\beta$.
\end{proposition}

\begin{proof}
The slope is $\widehat{\operatorname{Cov}}(x, y)/\widehat{\operatorname{Var}}(x)$, which converges to $\operatorname{Cov}(x^* + e, \beta x^* + u)/\operatorname{Var}(x^* + e)$. By independence the
numerator is $\beta\operatorname{Var}(x^*)$ and the denominator $\operatorname{Var}(x^*) + \operatorname{Var}(e)$.
\end{proof}

\begin{definition}[Attenuation bias]\label{def:b2:specification:attenuation}
\emph{Attenuation bias}\index{attenuation bias} is the shrinking of a coefficient toward zero caused by \omterm{def:b2:specification:classical-me}{classical
measurement error} in its regressor.
\end{definition}

Error in the outcome is different: it adds to $u$, so it raises the standard errors but leaves the slope unbiased.

\omcode{code/bachelor-2/32-specification/listings/attenuation.py}{1}{21}{The same noise added to the regressor, then to
the outcome.}\label{lst:b2:specification:attenuation}

\omoutput{code/bachelor-2/32-specification/outputs/attenuation.txt}{Noise in $x$ shrinks the slope to $\lambda\beta$; noise
in $y$ leaves it at 2.}\label{out:b2:specification:attenuation}

With several regressors the attenuation is stronger: what matters is the reliability of the part of $x$ not predicted by
the other regressors, and controls remove true variation while leaving the noise. Errors that are not classical (people
who under-report high incomes, for instance) can bias in either direction.

\section{Proxies and bad controls}\label{sec:b2:specification:controls}

\begin{definition}[Proxy variable]\label{def:b2:specification:proxy}
A \emph{proxy variable}\index{proxy variable} is an observed variable included in place of an unobserved one, because it
is correlated with it.
\end{definition}

A test score stands in for ability in a wage regression. If the score were a perfect measure, controlling for it would
remove the \omterm{def:b2:multiple-regression:ovb}{omitted-variable bias}; a noisy score removes part of it. The listing simulates it.

\omcode{code/bachelor-2/32-specification/listings/proxy.py}{1}{22}{Ability, a noisy test score, and the return to
schooling.}\label{lst:b2:specification:proxy}

\omoutput{code/bachelor-2/32-specification/outputs/proxy.txt}{The proxy removes more than half the bias, not all of
it.}\label{out:b2:specification:proxy}

% code: listings/proxy.py
The true return is 0.080; leaving ability out gives 0.105; the score brings it to 0.091. A proxy is a good control: it
is determined before the regressor of interest and is not changed by it. A control that the regressor itself changes is
not.

\begin{definition}[Bad control]\label{def:b2:specification:bad-control}
A \emph{bad control}\index{bad control} is a variable that is itself an outcome of the regressor of interest; holding it
fixed removes part of the effect being measured, and can create bias where there was none.
\end{definition}

Occupation in a regression of wages on schooling is the standard example: schooling changes occupation, so comparing
people with different schooling in the same occupation compares a clerk who studied long, and so is unusual among clerks,
perhaps unlucky or less able, with a clerk who did not. The 1975 earnings of \cref{ch:b2:multiple-regression} were a \omterm{def:b2:specification:bad-control}{bad control} of the same kind.

\section{Searching over specifications}\label{sec:b2:specification:search}

\begin{definition}[Specification search]\label{def:b2:specification:spec-search}
A \emph{specification search}\index{specification search} is the practice of estimating many specifications (other
regressors, forms, samples) and reporting those whose results look best; the reported tests then no longer have their
stated size.
\end{definition}

A 5 per cent test rejects a true null once in twenty. Try twenty unrelated regressors and keep the best, and the chance
that at least one is ``significant'' is $1 - 0.95^{20} = 0.642$.

\omcode{code/bachelor-2/32-specification/listings/search.py}{1}{20}{The best of twenty unrelated regressors, 5,000 times.}\label{lst:b2:specification:search}

\omoutput{code/bachelor-2/32-specification/outputs/search.txt}{The best of twenty is significant 64 per cent of the time;
the Bonferroni threshold brings it below 5 per cent.}\label{out:b2:specification:search}

% code: listings/search.py
The simulation gives 0.640. The Bonferroni correction divides the significance level by the number of tests, here
$0.05/20 = 0.0025$, a $|t|$ above about 3.1 with the listing's 98 degrees of freedom (3.02 in large samples); by the union bound the chance that any of the twenty passes is then at most
5 per cent (0.045 in the simulation). The correction works only if every test tried is counted, which a reader cannot
check; writing down the specification before seeing the data is the remedy that can be checked.

\section{Functional form and outliers}\label{sec:b2:specification:form}

A regression that is linear in the wrong variables misses curvature. One check adds the squares and cubes of the fitted
values to the regression and tests them jointly with an $F$ statistic: if they help predict $y$, the form is wrong
somewhere (this is the RESET test). It says that something is missing, not what. Outliers are another matter: an
observation far from the others in its regressors has high leverage and can move the fitted line on its own. Dropping it
because it is inconvenient is a \omterm{def:b2:specification:spec-search}{specification search}; the honest course is to report the result with and without it, and
to check that it is not a recording error.

\begin{model}{Classical measurement error}
\assumes The observed regressor is the true value plus an error with mean zero, independent of the true value and of the
regression error.
\predicts The slope converges to $\lambda\beta$, shrunk toward zero by the \omterm{def:b2:specification:reliability}{reliability ratio}; error in the outcome leaves
the slope unbiased.
\tested Where a variable is measured twice, by survey and by a record, the two can be compared: the \omterm{def:b2:specification:reliability}{reliability ratio}
of reported schooling or earnings can be estimated, and corrected slopes compared with uncorrected ones.
\breaks Errors are often not classical: people who earn more may under-report more, so the error is negatively
correlated with the true value; binary variables cannot have classical error at all, since a misreported 0 or 1 is
always wrong in a direction that depends on the true value.
\end{model}

\section{Exercises}\label{sec:b2:specification:exercises}

\begin{exercise}[$\star$]\label{exo:b2:specification:1}
Read the opening chart. What is the estimated slope when the error variance equals the variance of the true regressor?
What is the \omterm{def:b2:specification:reliability}{reliability ratio} then?
\end{exercise}

\begin{exercise}[$\star$]\label{exo:b2:specification:2}
Interpret the Engel elasticity of 0.856. If income rises by 10 per cent, by how much does food spending rise, and what
happens to the budget share?
\end{exercise}

\begin{exercise}[$\star$]\label{exo:b2:specification:3}
Why does the Brenland regression have four size dummies for five sizes? By how much, in per cent, does a household of
three spend more on food than a single person with the same income?
\end{exercise}

\begin{exercise}[$\star$]\label{exo:b2:specification:4}
A study reports that one of twenty diet variables significantly predicts a disease at the 5 per cent level, and says
nothing of the other nineteen. What should a reader conclude?
\end{exercise}

\begin{exercise}[$\star\star$]\label{exo:b2:specification:5}
In a regression of log wages on a dummy for union membership, the coefficient is 0.15. Give the approximate and the exact
percentage difference.
\end{exercise}

\begin{exercise}[$\star\star$]\label{exo:b2:specification:6}
Predict, then compute: what error variance makes the attenuated slope in the listing equal to 1.6? To 1?
\end{exercise}

\begin{exercise}[$\star\star$]\label{exo:b2:specification:7}
\textbf{(D)} With the Engel data, regress log food spending on centred log income and its square. Does the elasticity
change with income? Give it at the 10th and 90th percentiles of income.
\end{exercise}

\begin{exercise}[$\star\star$]\label{exo:b2:specification:8}
Change the proxy listing so that the test score has noise 0.3 instead of 0.7. Predict, then compute, the coefficient on
schooling.
\end{exercise}

\begin{exercise}[$\star\star$]\label{exo:b2:specification:9}
To evaluate a training programme on wages, a researcher controls for the worker's occupation after the programme. Why is
that a \omterm{def:b2:specification:bad-control}{bad control}? Which controls would be good?
\end{exercise}

\begin{exercise}[$\star\star\star$]\label{exo:b2:specification:10}
With $y = \beta_0 + \beta_1x_1^* + \beta_2x_2 + u$ and $x_1 = x_1^* + e$ classical, show using the \omterm{def:b2:multiple-regression:fwl}{Frisch--Waugh--Lovell
theorem} that the slope on $x_1$ converges to $\lambda_1\beta_1$, where $\lambda_1$ is the reliability of $x_1$ after
partialling out $x_2$, and that $\lambda_1 \leq \lambda$.
\end{exercise}

\begin{exercise}[$\star\star\star$]\label{exo:b2:specification:11}
Show that classical error in the outcome leaves the slope consistent, and give its effect on the variance of the slope.
\end{exercise}

\begin{exercise}[$\star\star\star$]\label{exo:b2:specification:12}
Show that if $k$ tests each have size $\alpha$, the probability that at least one rejects is at most $k\alpha$, and equals
$1 - (1 - \alpha)^k$ if they are independent. Compute both for $k = 20$, $\alpha = 0.05$ and for $\alpha = 0.0025$.
\end{exercise}

\section{Problem: Brenland's food survey}\label{sec:b2:specification:problem}

\begin{problem}[{Weekend problem --- Brenland's food survey: logs, dummies and interactions, reported income with a known
reliability, and a specification search}]
\label{pb:b2:specification:1}
Brenland's household survey records each household's food spending, disposable income, size and whether it owns its
home. A validation study, which compared reported incomes with tax records, found that the \omterm{def:b2:specification:reliability}{reliability ratio} of reported
log disposable income, after partialling out household size, is 0.8.

\medskip
\noindent\textbf{Part I --- Logs.}
\begin{enumerate}
  \item Regress log food on log disposable income. Give the elasticity and $R^2$.
  \item Is food a necessity or a luxury in Brenland?
  \item If income rises 10 per cent, by how much does food spending rise?
  \item What does the budget share of food do as income rises?
  \item Why use logs rather than levels here?
\end{enumerate}

\medskip
\noindent\textbf{Part II --- Dummies and interactions.}
\begin{enumerate}[resume]
  \item Add a dummy for each household size from 2 to 5 and an owner dummy. Give the elasticity and $R^2$.
  \item Why did the elasticity fall?
  \item Give the size coefficients and the percentage difference for a household of five.
  \item Why is there no dummy for size 1?
  \item Interpret the owner coefficient.
  \item Add the interaction of owner and log income. Give the elasticities of owners and tenants.
  \item Is the difference significant at 5 per cent?
\end{enumerate}

\medskip
\noindent\textbf{Part III --- Reported income with a known reliability.}
\begin{enumerate}[resume]
  \item If reported income has classical error, which way is the elasticity of Part II biased?
  \item Correct it with the \omterm{def:b2:specification:reliability}{reliability ratio}.
  \item Why is the reliability after partialling out size the right one?
  \item Would the correction be the same without the size dummies?
  \item Would error in reported food spending call for a correction?
  \item What if high earners under-report most?
\end{enumerate}

\medskip
\noindent\textbf{Part IV --- A specification search and an honest threshold.}
\begin{enumerate}[resume]
  \item An analyst tries twenty other household characteristics, one at a time, as extra regressors. If none matters, what
  is the chance that at least one is significant at 5 per cent?
  \item Give the Bonferroni threshold for the $p$-value and the matching $|t|$.
  \item Would the interaction of Part II pass it?
  \item What should the analyst write down before looking at the data?
  \item Why is household size not part of the search?
  \item Could the analyst choose the twenty characteristics after seeing which look promising?
  \item State the result: the corrected food--income elasticity and the Bonferroni threshold for twenty tries.
\end{enumerate}
\end{problem}
